\chapter{Player Classification Rules}
\label{app:player_classification}

\section*{B.1\quad Minimum Data Threshold}

A player is assigned the label \texttt{INSUFFICIENT\_DATA} and excluded from substantive classification if either of the following conditions holds:

\begin{itemize}
    \item The player has appeared in fewer than \textbf{15 matches} in the dataset.
    \item The player has accumulated fewer than \textbf{250 total runs} \textit{and} fewer than \textbf{10 total wickets}.
\end{itemize}

Players meeting the minimum threshold are then assessed by the ordered rule set in Section~B.2. The first rule that matches determines the final classification; subsequent rules are not evaluated.

\section*{B.2\quad Ordered Classification Rules}

\begin{table}[ht]
\centering
\caption{Player classification rules applied in priority order}
\label{tab:classification_rules_full}
\begin{tabular}{p{4cm} p{9.5cm}}
\hline
\textbf{Classification} & \textbf{Conditions (all must hold)} \\
\hline
\textbf{BOWLER} &
  Overs bowled per match $\geq 3$ \newline
  Economy rate $\leq 8.2$ runs per over \newline
  Bowling average $\leq 26$ runs per wicket \newline
  Batting average $< 20$ runs per dismissal \\
\hline
\textbf{ALL\_ROUNDER} \newline (bowling-dominant) &
  Overs bowled per match $\geq 2$ \newline
  Total wickets $\geq 25$ \newline
  Batting average $\geq 20$ runs per dismissal \newline
  Batting strike rate $\geq 120$ \newline
  Economy rate $\leq 8.8$ runs per over \\
\hline
\textbf{ALL\_ROUNDER} \newline (batting-dominant) &
  Overs bowled per match $\geq 1$ \newline
  Total wickets $\geq 20$ \newline
  Batting average $\geq 25$ runs per dismissal \newline
  Batting strike rate $\geq 130$ \\
\hline
\textbf{BATTER} &
  Batting average $\geq 28$ runs per dismissal \newline
  Batting strike rate $\geq 125$ \newline
  Overs bowled per match $< 0.5$ \\
\hline
\textbf{UTILITY\_PLAYER} &
  All remaining players who meet the minimum data threshold and do not match any rule above. \\
\hline
\end{tabular}
\end{table}

\noindent Both ALL\_ROUNDER sub-types receive the same \texttt{ALL\_ROUNDER} label and the same RBSS weight (0.80). The distinction between bowling-dominant and batting-dominant all-rounders is made only for internal rule ordering and does not appear in the prompt.

\section*{B.3\quad RBSS Weight Design and Rationale}

The RBSS weight for each classification was selected to reflect expected batting contribution in T20 cricket, not bowling quality. The weights are designed to be \textit{stable}, \textit{simple}, \textit{rounded}, and \textit{consistent} across the full dataset, as language models interpret smooth and interpretable numeric values more reliably than arbitrary decimals.

\begin{table}[ht]
\centering
\caption{RBSS weights assigned to each player classification}
\label{tab:rbss_weights_full}
\begin{tabular}{lcp{7.5cm}}
\hline
\textbf{Classification} & \textbf{Weight ($w$)} & \textbf{Design rationale} \\
\hline
BATTER             & 1.00 & Baseline reference unit. Specialist batters are the primary run-scoring resource; all other weights are calibrated relative to this value. \\
ALL\_ROUNDER       & 0.80 & T20 all-rounders typically bat at positions 5--7 with high strike rates but lower consistency than top-order batters. Bowling proficiency partially substitutes for batting depth, justifying a modest discount from 1.00. \\
UTILITY\_PLAYER   & 0.65 & Floaters and situational players with mixed skill sets and inconsistent output. Positioned between reliable batting contributors and tail-end bowlers. \\
INSUFFICIENT\_DATA & 0.45 & Players with limited match history. A low weight (e.g.\ 0.30) would systematically undervalue unknown young batters; a high weight (e.g.\ 0.80) would be overly optimistic. The value 0.45 is a conservative mid-point that minimises expected error under uncertainty. \\
BOWLER             & 0.30 & Specialist bowlers bat in the lower order (positions 8--11). The weight is non-zero to capture the occasional late-innings cameo hitting that is common in T20 cricket, but is kept low to reflect the generally limited batting expectation. \\
\hline
\end{tabular}
\end{table}

\subsection*{B.3.1\quad Formula}

The RBSS is computed at the end of each two-over section as:
\begin{equation}
    \text{RBSS} = \sum_{p \,\in\, \mathcal{R}} w\bigl(\text{class}(p)\bigr)
\end{equation}

where $\mathcal{R}$ is the set of players from the batting team's eleven who have not yet been dismissed at the end of that section, and $w(\cdot)$ maps each classification to its weight. The score is rounded to two decimal places.

\subsection*{B.3.2\quad Worked Example}

Consider a match state at the end of over 12 in which the following players remain undismissed:

\begin{center}
\begin{tabular}{ll}
\hline
\textbf{Player} & \textbf{Classification} \\
\hline
Player A & BATTER \\
Player B & ALL\_ROUNDER \\
Player C & ALL\_ROUNDER \\
Player D & UTILITY\_PLAYER \\
Player E & BOWLER \\
Player F & BOWLER \\
Player G & BOWLER \\
\hline
\end{tabular}
\end{center}

\vspace{0.5em}
The RBSS at this point is:

\begin{equation}
    \text{RBSS} = (1 \times 1.00) + (2 \times 0.80) + (1 \times 0.65) + (3 \times 0.30)
    = 1.00 + 1.60 + 0.65 + 0.90 = 4.15
\end{equation}

This is encoded in the natural-language prompt as \texttt{Remaining Batting Strength Score: 4.15}. The model can infer from this value that the batting side has limited remaining depth --- a team of seven pure batters would yield an RBSS of 7.00, while seven pure bowlers would yield only 2.10.

\subsection*{B.3.3\quad Boundary Cases}

A full eleven of \texttt{INSUFFICIENT\_DATA} players at the start of an innings yields a maximum RBSS of $11 \times 0.45 = 4.95$. A team composed entirely of specialist batters would yield $11 \times 1.00 = 11.00$. In practice, RBSS values at the start of an innings typically fall in the range 6.0--8.5, decreasing as wickets fall and the lower-order bowlers approach the crease.

\section*{B.4\quad Implementation Reference}

The classification logic is implemented in \texttt{forecaster/player\_classification.py}, function \texttt{classify\_t20i\_player}. Player career statistics used as inputs to the classification rules are aggregated across all qualifying first-innings match rows in the dataset and stored in the SQLite table \texttt{first\_innings\_player\_stats}.
